\section{Итоги}
\label{sec:summary}

\textbf{Равновесный сдвиг.} Периодичность синхротронного движения даёт точное условие $\langle\tau\rangle=0$, откуда
\[
\delta_\eq=-\frac{1}{\eta_0}\Biggl(\frac{\eta_1\delta_m^2}{2}+\Lambda\Biggr)
\]
(формула~\eqref{eq:sol-deq}). Знак и делитель $\eta_0$ не зависят от договорённости о знаке фазы и от приближения $\cos\varphi\approx 1$; метод Боголюбова--Митропольского нужен только для связи $\langle\delta^2\rangle=\delta_m^2/2$ с амплитудой первой гармоники. Следующий член по амплитуде есть $-(19/72)\,\eta_1^3\delta_m^4/\eta_0^3$~\eqref{eq:sol-a4}.

\textbf{Формула (10) работы~\cite{Senichev2013}} воспроизводится с двумя поправками. Множитель $\gamma_s^2/(\gamma_s^2\alpha_0-1)$ равен $+1/\eta_0$, а должен стоять $-1/\eta_0$ (раздел~\ref{sec:mag}). В скобке при $\delta_m^2/2$ не хватает слагаемого $3\beta_s^2/(2\gamma_s^2)$: полный коэффициент есть $\eta_1$ из~\eqref{eq:mag-eta1}, а обрыв $\Delta v/v=\delta/\gamma^2$ теряет вторую производную скорости.

\textbf{Поправки к решению.} Вторая гармоника равна $\eta_1\delta_m^2/(6\eta_0)\cos 2\psi$; в $\delta_\eq$ на порядке $\delta_m^2$ она не входит. Частота синхротронных колебаний~\eqref{eq:sol-omega-sum}
\[
\frac{\Omega}{\omega_s}=1-\frac{\eta_1\Lambda}{\eta_0^2}-\frac{5}{12}\Biggl(\frac{\eta_1\delta_m}{\eta_0}\Biggr)^2-\frac{\varphi_m^2}{16}
\]
содержит все три поправки; поправка конечной фазы $-\varphi_m^2/16$ в $\delta_\eq$ входит лишь порядком $\varphi_m^4\delta_m^2$.

\textbf{Член $(\Delta L/L)_\beta$ в формуле (6)} в магнитном кольце стоит с коэффициентом $1$, потому что скорость вдоль траектории постоянна и отношение времён равно отношению длин (раздел~\ref{sec:beta}). Общая формула $\Lambda$ через $\gamma_{x,y}$ Твисса~\eqref{eq:beta-Lambda} переходит в формулу (11) в гладком пределе при $\varepsilon=2J$.

\textbf{Определения Нга.} Вывод использует делитель $T_s$, как (18.10) в~\cite{Ng2006}. Определения (18.20) и (18.21) дали бы $\eta_1-\eta_0^2$ и $2(\eta_1-\eta_0^2)$ и сдвинули бы $\delta_\eq$~\eqref{eq:ng-deq}. Ни наш вывод, ни формулы (3)--(5) работы~\cite{Senichev2013} этой проблемы не имеют. При чтении COSY и MAD-X нужно помнить множитель $2$ между коэффициентом ряда и производной и пересчёт от $\delta_K$ или $\mathrm{PT}$ к $\delta$.

\textbf{Электростатика.} Сохраняется полная энергия $\gamma mc^2+q\Phi$, и медленная переменная $\delta$ определена через неё; уравнения зазора и фазы не меняются, меняется только $\tau(\delta,J)$. На смещённой траектории частица движется с другой скоростью, поэтому задержка равна удлинению пути плюс вклад изменённой скорости. Линейный коэффициент $\eta_0=\alpha_0+\alpha_0^E/\gamma_s^2-1/\gamma_s^2$~\eqref{eq:es-eta0} (в фильтре Вина $\alpha_0^E\neq 0$ при нулевой кривизне); квадратичный~\eqref{eq:es-eta1} содержит $\langle D_2/\rho\rangle$ и $\langle(D/\rho)^2\rangle$, которые через $\alpha_0$, $\alpha_1$ не выражаются; бетатронный вклад~\eqref{eq:es-lambda} к одному множителю перед $\Delta L/C_s$ не сводится. В гладком кольце с цилиндрическим дефлектором $\Lambda^{\mathrm{es}}=0$ при ненулевом удлинении пути, со сферическим $\Lambda^{\mathrm{es}}=-\beta_s^2\langle(x/\rho)^2\rangle$.

\textbf{Формула (25).} Множитель $F_v=1/(1-\alpha_0^E)$ есть отношение приращения скорости вне поля к приращению на орбите~\eqref{eq:f25-fv}; при нём линейный член (25) точен. Квадратичный член верен для цилиндрического дефлектора и неверен для сферического: пропущен аналог $3\beta_s^2/(2\gamma_s^2)$, а степень $F_v^2$ при $\alpha_0/\gamma^2$ из исходной логики не следует. Коэффициент $1$ перед $(\Delta L/L)_\beta$ в электростатике неверен.
